% GENERATED by docs/paper/scripts/tables.py from coco_lab/coco_lab/benchmarks/*.v1.json
\begin{table}[t]
\centering
\small
\caption{COCO-Bench v1. Families: P planning, C control, L localisation, T full task. Tiers easy/medium/hard; held-out scenarios drawn per evaluation.}
\label{tab:suites}
\setlength{\tabcolsep}{3pt}
\begin{tabular}{lcrcrl}
\toprule
suite & fam. & $n$ & tiers & held-out & world \\
\midrule
\texttt{arena-nav-50} & PCL & 50 & 17/17/16 & 15 & \texttt{arena} \\
\texttt{ctrl-corners-30} & C & 30 & 10/10/10 & 10 & \texttt{depot} \\
\texttt{ctrl-corridor-30} & C & 30 & 10/10/10 & 10 & \texttt{depot} \\
\texttt{ctrl-dynamic-30} & C & 30 & 10/10/10 & 10 & \texttt{depot} \\
\texttt{depot-clutter-50} & PC & 50 & 17/17/16 & 15 & \texttt{depot\_clutter} \\
\texttt{depot-long-50} & PC & 50 & 17/17/16 & 15 & \texttt{depot} \\
\texttt{depot-narrow-30} & PC & 30 & 10/10/10 & 10 & \texttt{depot} \\
\texttt{depot-nav-100} & PCL & 100 & 34/33/33 & 20 & \texttt{depot} \\
\texttt{loc-global-20} & L & 20 & 7/7/6 & 6 & \texttt{depot} \\
\texttt{loc-kidnap-20} & L & 20 & 7/7/6 & 6 & \texttt{depot} \\
\texttt{loc-track-30} & L & 30 & 10/10/10 & 10 & \texttt{depot} \\
\texttt{task-fetch-10} & T & 10 & 4/3/3 & 0 & \texttt{depot+depot\_colours} \\
\bottomrule
\end{tabular}
\end{table}


\begin{figure*}[t]
\centering
\includegraphics[width=\textwidth]{figures/tiers.pdf}
\caption{Planning scenarios of four \bench{} v1 suites by the reference path's length and bends,
coloured by tier (generated from the frozen suite files).}
\label{fig:tiers}
\end{figure*}

\textbf{Suites.} \bench{} v1 (\cref{tab:suites}) has \BenchSuites{} frozen suites with
\BenchScenarios{} scenarios on the original COCO arena and on a \DepotWidthM{}~m by \DepotHeightM{}~m warehouse, the
Depot, with \DepotZones{} zones, rack aisles, a workshop maze and \DepotActors{} scheduled forklifts and walkers. \BenchCarried{}
global-planning suites carry over the earlier v0 suites with identical scenarios (v1 adds
features and tiers only); \texttt{depot-long-50} keeps routes whose reference path is long and
crosses several zones; \texttt{depot-clutter-50} runs on a variant of the Depot with
\ClutterPallets{} extra pallets taken from the world's own kit, placed by a seeded generator that
keeps every carried-over start and goal connected, so a scenario looks exactly as cluttered as it
is in the model and in Gazebo. The local-control suites select long narrow passages, routes with
many bends, and routes that pass the actors' lanes; in \texttt{ctrl-dynamic-30} each start and
goal appears twice, once with the actors stopping for the robot and once with them keeping their
schedule, a scenario rule rather than physics. The localisation suites test tracking from a
known start, global localisation from an unknown start (filters that can represent it), and the
kidnapped robot. The full-task suite is the committed fetch rounds; we did not invent new ones.

\textbf{Features and tiers.} Each scenario is described by features of one reference plan (A* on
the default costmap): path length, the narrowest clearance, the longest narrow stretch, bends of
\BendDeg{}$^\circ$ or more, zones crossed and exposure to actor lanes. Within a suite, each tier
feature becomes a percentile rank, the ranks are averaged into a difficulty score, and the
scenarios are cut into thirds (\cref{fig:tiers}). Results are reported per tier.

\textbf{Metrics.} Planning: success, path length (and its ratio to the grid-optimal length),
path cost, the minimum and mean clearance and expansions. Control: success under Nav2-like
goal and progress checkers, time, the distance driven by the true pose, clearance and contacts.
Localisation: success (reaching the goal with final error below thresholds), final, mean and
maximum error, convergence and recovery time. Success rates come with Wilson intervals
\cite{wilson1927}; continuous metrics with medians and quartiles. Wall-clock compute time is
reported but never compared across machines and never enters a hash.

\textbf{Held-out scenarios.} A leaderboard entry is also run on scenarios it has never seen:
\(\mathrm{seed} = \mathrm{HMAC}(\mathit{secret},\ \mathit{suite} \,\|\, \mathit{head})\), where
\(\mathit{head}\) is the submission's commit, so every push gets a fresh set. The trusted
workflow generates the scenario file before any untrusted code runs and passes only that file
into the sandbox; the secret never enters the container, the container's output is never
published, and only rounded aggregates leave. \Cref{sec:limits} states what this does not
prevent.

\textbf{Datasheets.} Each suite ships a generated datasheet: purpose, composition, generation
procedure and seeds, conditions, the features per tier, metric definitions with formulas,
held-out policy, known biases and limitations, versioning and licence.
